\section{Introduction}
\label{sec:introduction}
The rapid growth of electric mobility is shifting charging-infrastructure planning from static placement rules toward data-driven decision support \cite{iea2023evoutlook}. Bidirectional charging adds a further dimension, since vehicles are no longer only energy consumers but can also feed power back to the grid, a concept known as \ac{V2G}, which turns charging points into distributed flexibility assets \cite{kempton2005v2g}. Deciding where to place such stations therefore requires understanding where and when charging demand actually concentrates.

Because electric vehicles still represent a small share of the fleet, there is little real-world data on where charging stations should be placed and how intensively they are used. At the same time, uncoordinated charging concentrates electrical load in space and time and can stress local distribution grids \cite{clementnyns2010impact}, which makes the spatial and temporal structure of demand a central planning question. Existing siting studies, such as \cite{he2013siting}, however, often rely on data or tooling that is not openly reproducible. We therefore ask: \emph{how can positions for optimal charging station placement (candidate locations) be ranked from openly available data, based on their accessibility and distribution of charging demand over space and time?}

To address this, we build a simulation workflow on openly available data. Road-network and land-use information are taken from \ac{OSM}, and traffic is simulated microscopically in \ac{SUMO}, controlled through its Python interface \ac{TraCI}. From  \acsp{POI}  we derive travel demand, and model charging as a \emph{spatiotemporal} process, that is, demand that varies simultaneously by location and time of day. The model distinguishes public charging stations from private home wallboxes with vehicle-specific access control. A custom web interface, which is based on \ac{SUMO}'s web interface, lets users select an area of interest on a map and generate ready-to-run scenarios directly, lowering the barrier to reproducing and extending the analysis. Candidate sites are then assessed in terms of accessibility, station utilization, and the resulting local charging load.

The scope of this study is the planning and evaluation stage. We focus on the spatial and temporal distribution of charging demand under realistic mobility patterns. The distribution grid is represented by a synthetic network whose power flow bounds the power available at each station, so that the charging we report stays within its voltage and loading limits. Analysing that grid is not the subject of this study. Peak load, hosting capacity and reinforcement planning are left to future work, as are battery degradation and market or pricing mechanisms.